

\chapter{Modelo C4 de Arquitectura}

El modelo C4 (Context, Containers, Components, Code) propuesto por Simon Brown ofrece una manera abstracta de comunicar la arquitectura de un sistema de software a distintos niveles de detalle. Este capítulo presenta los tres primeros niveles para el sistema FisioManager.

\section{Nivel 1: Diagrama de Contexto}

El diagrama de contexto muestra a FisioManager como una caja azul y describe sus interacciones con usuarios y sistemas externos.

\begin{figure}[H]
  \centering
  \includegraphics[width=0.85\textwidth,keepaspectratio]{FisioManager_C4_Contexto.png}
  \caption{Diagrama de Contexto del Sistema (C4 Nivel 1).}
  \label{fig:c4-contexto}
\end{figure}

\subsection{Actores y sistemas externos}

\begin{itemize}
  \item \actor{Paciente:} interactúa con FisioManager vía web responsiva para visualizar rutinas, ejecutar ejercicios y registrar feedback por voz.
  \item \actor{Terapeuta:} gestiona pacientes, crea rutinas, asigna y supervisa progreso.
  \item \actor{Administrador:} actor único que supervisa indicadores agregados de uso y desempeño de la clínica y, a la vez, gestiona usuarios, accesos y monitoreo técnico de la plataforma.
  \item \actor{Ollama (LLM local):} genera el resumen clínico del feedback del paciente.
  \item \actor{Almacenamiento de Objetos (MinIO/S3):} persiste avatares, documentos clínicos y audios.
  \item \actor{Proveedor de Email (SMTP):} envía notificaciones y recordatorios por correo.
\end{itemize}

La transcripción de voz (Whisper) se ejecuta dentro del propio sistema, por lo que no constituye un sistema externo.

\section{Nivel 2: Diagrama de Contenedores}

El diagrama de contenedores descompone FisioManager en sus aplicaciones, servicios y almacenes de datos. Cada contenedor representa una unidad desplegable independientemente.

\begin{figure}[H]
  \centering
  \includegraphics[width=0.95\textwidth,keepaspectratio]{FisioManager_C4_Contenedores.png}
  \caption{Diagrama de Contenedores (C4 Nivel 2).}
  \label{fig:c4-contenedores}
\end{figure}

\subsection{Contenedores principales}

\begin{longtable}{L{3.5cm} L{3cm} L{8cm}}
  \toprule
  \rowcolor{fmTableHeader}
  \textbf{\color{fmAccent}Contenedor} &
  \textbf{\color{fmAccent}Tecnología} &
  \textbf{\color{fmAccent}Responsabilidad} \\
  \midrule
  \endhead

  \rowcolor{fmTableAltRow}
  \actor{Reverse-proxy / TLS} &
  Caddy / nginx &
  Termina HTTPS, sirve los archivos estáticos de la SPA y enruta \texttt{/api}, \texttt{/uploads} y \texttt{/socket.io} hacia la API. \\

  \actor{Aplicación Web SPA} &
  React + TypeScript + Vite (PWA) &
  Interfaz de usuario única para los tres actores. Captura audio en el navegador, soporta i18n, modo claro/oscuro y uso offline. \\

  \rowcolor{fmTableAltRow}
  \actor{API REST + WebSocket} &
  Node.js + Express + Socket.IO &
  Monolito modular (rutas → controladores → servicios → repositorios). Autenticación JWT, autorización por rol, validación con Zod, rate limiting, auditoría y acceso a datos con Prisma. \\

  \actor{Transcripción de voz} &
  Python + Whisper &
  Convierte el audio del feedback a texto; se ejecuta como subproceso dentro de la imagen de la API. \\

  \rowcolor{fmTableAltRow}
  \actor{Base de Datos} &
  PostgreSQL 16 (Prisma) &
  Persistencia transaccional principal: usuarios, historial clínico, rutinas, sesiones, feedback, citas, mensajes y auditoría. \\

  \actor{Cache / PubSub} &
  Redis &
  Adapter de Socket.IO para propagar mensajes y notificaciones entre varias instancias de la API (escalado horizontal). \\

  \rowcolor{fmTableAltRow}
  \actor{Object Storage} &
  MinIO / S3 &
  Almacena avatares y documentos clínicos cuando se activa el modo de almacenamiento de objetos. \\

  \actor{Respaldos de BD} &
  pg\_dump programado &
  Genera copias de seguridad periódicas de PostgreSQL con retención configurable. \\

  \bottomrule
\end{longtable}

Ollama (LLM local) y el proveedor SMTP se consideran sistemas externos a la frontera de despliegue.

\section{Nivel 3: Diagrama de Componentes}

El nivel 3 descompone la \textbf{API} en sus componentes internos, tomando como ejemplo el flujo de \textbf{Feedback}, por ser el de mayor complejidad de orquestación (audio, transcripción y resumen con IA, alertas). Todos los módulos siguen el mismo patrón en capas: \textit{rutas → controladores → servicios → repositorios}.

\begin{figure}[H]
  \centering
  \includegraphics[width=0.95\textwidth,keepaspectratio]{FisioManager_C4_Componentes.png}
  \caption{Componentes del Servicio de Feedback (C4 Nivel 3).}
  \label{fig:c4-componentes}
\end{figure}

\subsection{Componentes del flujo de Feedback}

\begin{itemize}
  \item \textbf{Middlewares:} \texttt{auth} (verificación JWT), \texttt{role} (autorización por rol), \texttt{validate} (esquemas Zod) y \texttt{rate-limit} se aplican antes de cada ruta.
  \item \textbf{feedback.routes:} define los endpoints \texttt{POST /feedback} y \texttt{POST /feedback/audio} sobre un Router de Express.
  \item \textbf{feedback.controller:} adapta la petición/respuesta (\texttt{catchAsync}, helpers \texttt{ok}/\texttt{created}) y delega en los servicios.
  \item \textbf{feedback.service:} implementa las reglas de negocio del registro de feedback y dispara la alerta de dolor alto.
  \item \textbf{ai.service:} ejecuta \texttt{transcribeAudio} con Whisper y \texttt{generateClinicalSummary} con Ollama.
  \item \textbf{storage (lib):} guarda el audio en disco local o en MinIO/S3 según la configuración.
  \item \textbf{notification.service:} crea la notificación in-app, envía el correo y empuja el evento por Socket.IO.
  \item \textbf{feedback.repository / painPoint.repository:} encapsulan el acceso a datos mediante Prisma Client.
\end{itemize}